% ============================================================
%  Section: System Model and Problem Formulation (Single-UAV)
%  (可直接粘贴到 Overleaf 主文档中; 需要 amsmath, amssymb 宏包)
% ============================================================

\section{SYSTEM MODEL AND PROBLEM FORMULATION}
\label{sec:system_model}

% ==============================================================
\subsection{Problem Statement}
\label{subsec:problem_statement}

We consider the target coverage problem in a three-dimensional environment
composed of a UAV, $M$~mobile target zones
$\mathcal{G} = \{G_1, \dots, G_M\}$, and $K$~dynamic obstacle zones
$\mathcal{T} = \{T_1, \dots, T_K\}$, collectively referred to as entities
$\mathcal{E} = \{\mathcal{G}, \mathcal{T}\}$.
At the beginning of each episode, the positions of all entities are randomly
initialised within a bounded task region.
The UAV has a fixed observation radius~$R_{\mathrm{obs}}$ and a coverage
radius~$R_{\mathrm{cov}}$, where $R_{\mathrm{obs}}$ defines the maximum
sensing range for detecting entities, and $R_{\mathrm{cov}}$ defines the
proximity required to cover a target.
The system objective is to \emph{maximise the cumulative target zone coverage within
a finite time horizon~$T_{\max}$ while ensuring collision avoidance with all
obstacle zones}.

% ==============================================================
\subsection{System Model}
\label{subsec:system_model}

We consider a bounded three-dimensional task region
$\mathcal{W} = [-L/2, L/2]^{2} \times [0, H_{\mathrm{w}}]$, where $L$
denotes the side length of the square operating area and $H_{\mathrm{w}}$ is
the height of the enclosing rigid boundaries.
The environment consists of the UAV, $M$~targets, and $K$~obstacle zones.
At time step~$t$, the UAV has position
$\mathbf{p}^{(t)} = (x^{(t)}, y^{(t)}, z^{(t)})^{\top}$ in the world frame,
velocity $\mathbf{v}^{(t)} = (\dot{x}^{(t)}, \dot{y}^{(t)},
\dot{z}^{(t)})^{\top} \in \mathbb{R}^3$, and its speed is bounded by
$v_{\max}$, i.e.\ $\lVert\mathbf{v}^{(t)}\rVert \le v_{\max}$.

\subsubsection{Quadrotor Dynamics}

The UAV is modelled as a quadrotor with mass~$m$, moment of inertia matrix
$\mathbf{J} \in \mathbb{R}^{3 \times 3}$, body radius~$r_{\mathrm{uav}}$, and
thrust coefficient~$k_f$.
The orientation is described by the Euler angles
$\boldsymbol{\Phi} = (\phi, \theta, \psi)^{\top}$, representing roll, pitch,
and yaw, respectively.
The four rotors ($l = 1,\dots,4$) rotate at angular speeds~$\omega_l$ and
produce thrusts $f_l = k_f\,\omega_l^2$.
The translational and rotational dynamics in the world frame are
\begin{align}
  m\,\ddot{\mathbf{p}} &=
    \mathbf{R}(\boldsymbol{\Phi})\,
    \begin{pmatrix}0\\0\\\sum_{l=1}^{4} f_l\end{pmatrix}
    - mg\,\mathbf{e}_3,
  \label{eq:trans_dyn} \\
  \mathbf{J}\,\dot{\boldsymbol{\omega}} &= \boldsymbol{\tau}
    - \boldsymbol{\omega} \times \mathbf{J}\,\boldsymbol{\omega},
  \label{eq:rot_dyn}
\end{align}
where $\mathbf{R}(\boldsymbol{\Phi}) \in SO(3)$ is the rotation matrix from
the body frame to the world frame, $g$ is the gravitational acceleration,
$\mathbf{e}_3 = (0,0,1)^{\top}$ is the unit vector along the world $z$-axis,
$\boldsymbol{\omega} \in \mathbb{R}^{3}$ is the body angular velocity, and
$\boldsymbol{\tau} = (\tau_\phi, \tau_\theta, \tau_\psi)^{\top}$ is the
body-frame torque vector derived from the rotor thrusts through the allocation
matrix.

\subsubsection{Cascaded PID Controller}

A cascaded PID controller maps the high-level velocity command to the rotor
speeds through three stages: a position loop, an attitude loop, and a motor
mixer.

\paragraph{Velocity command.}
Given the action $\mathbf{a} = (a^{x}, a^{y})^{\top} \in [-1,1]^{2}$
output by the policy, the desired velocity is
\begin{equation}
  \mathbf{v}^{\mathrm{cmd}} = v_{\max}\,
  \frac{\mathbf{a}}{\lVert\mathbf{a}\rVert},
  \label{eq:vel_cmd}
\end{equation}
where $v_{\max}$ is the maximum flight speed.
Only the $xy$-components are commanded; the altitude is held constant at
$z_0 = H_{\mathrm{w}}/2$.
The reference position is updated each control step as
$\mathbf{p}^{\mathrm{ref}} = \mathbf{p}^{(t)} +
\mathbf{v}^{\mathrm{cmd}} \cdot \Delta t$,
where $\Delta t$ is the control period.

\paragraph{Position loop.}
A PID controller computes the desired thrust vector from position and velocity
errors:
\begin{equation}
  \mathbf{F}^{\mathrm{des}} =
    \mathbf{K}_p^{\mathrm{pos}}\,\mathbf{e}_{p}
    + \mathbf{K}_i^{\mathrm{pos}}\!\int\!\mathbf{e}_{p}\,\mathrm{d}t
    + \mathbf{K}_d^{\mathrm{pos}}\,\mathbf{e}_{v}
    + mg\,\mathbf{e}_3,
  \label{eq:pid_pos}
\end{equation}
where $\mathbf{e}_{p} = \mathbf{p}^{\mathrm{ref}} - \mathbf{p}$ is the
position error, $\mathbf{e}_{v} = \mathbf{v}^{\mathrm{cmd}} - \mathbf{v}$ is
the velocity error, and
$\mathbf{K}_p^{\mathrm{pos}}, \mathbf{K}_i^{\mathrm{pos}},
\mathbf{K}_d^{\mathrm{pos}} \in \mathbb{R}^{3 \times 3}$ are the
proportional, integral, and derivative gain matrices of the position loop.
The scalar thrust is obtained by projecting $\mathbf{F}^{\mathrm{des}}$ onto
the current body $z$-axis, and the desired attitude
$\boldsymbol{\Phi}^{\mathrm{des}}$ is extracted from the direction
of~$\mathbf{F}^{\mathrm{des}}$.

\paragraph{Attitude loop.}
A second PID controller produces the desired body-frame torque:
\begin{equation}
  \boldsymbol{\tau}^{\mathrm{des}} =
    \mathbf{K}_p^{\mathrm{att}}\,\mathbf{e}_{\Phi}
    + \mathbf{K}_i^{\mathrm{att}}\!\int\!\mathbf{e}_{\Phi}\,\mathrm{d}t
    + \mathbf{K}_d^{\mathrm{att}}\,\mathbf{e}_{\omega},
  \label{eq:pid_att}
\end{equation}
where $\mathbf{e}_{\Phi}$ is the rotation error between the current attitude
$\boldsymbol{\Phi}$ and the desired attitude
$\boldsymbol{\Phi}^{\mathrm{des}}$,
$\mathbf{e}_{\omega}$ is the angular rate error, and
$\mathbf{K}_p^{\mathrm{att}}, \mathbf{K}_i^{\mathrm{att}},
\mathbf{K}_d^{\mathrm{att}} \in \mathbb{R}^{3 \times 3}$ are the attitude PID
gain matrices.

\paragraph{Motor mixer.}
The scalar thrust and desired torques are combined through a mixer matrix
$\mathbf{M} \in \mathbb{R}^{4\times3}$ to obtain the four motor commands:
\begin{equation}
  \boldsymbol{\omega}^{\mathrm{rpm}} =
    \alpha_{\mathrm{pwm}}\,\mathrm{clip}\!\bigl(
      T_{\mathrm{base}}\,\mathbf{1}_4
      + \mathbf{M}\,\boldsymbol{\tau}^{\mathrm{des}},\;
      \omega_{\min},\, \omega_{\max}
    \bigr)
    + \beta_{\mathrm{pwm}},
  \label{eq:mixer}
\end{equation}
where $T_{\mathrm{base}}$ is the hover thrust command,
$\mathbf{1}_4 = (1,1,1,1)^{\top}$ distributes it equally to four motors,
$\mathrm{clip}(\cdot,\, \omega_{\min},\, \omega_{\max})$ clamps the PWM signal
within the hardware limits $[\omega_{\min},\, \omega_{\max}]$, and
$\alpha_{\mathrm{pwm}}, \beta_{\mathrm{pwm}}$ are the affine coefficients of
the PWM-to-RPM conversion.

\subsubsection{Observation and Coverage Indicators}

The UAV monitors entities within its sensing range.
Let $\mathbf{p}_j^{(t)}$ denote the position of entity~$E_j$ at time step~$t$.
The observation indicator function is defined as
\begin{equation}
  I_{\mathrm{obs}}(E_j, t) =
  \begin{cases}
    1 & \text{if } \lVert\mathbf{p}^{(t)} - \mathbf{p}_j^{(t)}\rVert
        \le R_{\mathrm{obs}}, \\
    0 & \text{otherwise},
  \end{cases}
  \label{eq:obs_indicator}
\end{equation}
where $I_{\mathrm{obs}}(E_j, t) = 1$ indicates that entity~$E_j$ is detected
by the UAV at time~$t$.
Let $\mathbf{p}_m^{(t)}$ denote the position of target~$G_m$.
The coverage indicator function is
\begin{equation}
  I_{\mathrm{cov}}(G_m, t) =
  \begin{cases}
    1 & \text{if } \lVert\mathbf{p}^{(t)} - \mathbf{p}_m^{(t)}\rVert
        \le R_{\mathrm{cov}}, \\
    0 & \text{otherwise},
  \end{cases}
  \label{eq:cov_indicator}
\end{equation}
where $I_{\mathrm{cov}}(G_m, t) = 1$ indicates that target~$G_m$ is
successfully covered.
The objective is to maximise the cumulative coverage over the episode
horizon~$T$:
\begin{equation}
  J = \sum_{t=0}^{T} \sum_{m=1}^{M} I_{\mathrm{cov}}(G_m, t).
  \label{eq:coverage_obj}
\end{equation}
Upon coverage, the target is immediately respawned at a random collision-free
location to keep the task complexity constant.

\subsubsection{Dynamic Entities}

Each target zone $G_m$ is modelled as a cylindrical zone of radius~$r_{\mathrm{tgt}}$.
The target zone moves on the $xy$-plane at a constant speed~$v_{\mathrm{tgt}}$;
its heading is re-sampled uniformly from $\mathcal{U}(0, 2\pi)$ every
$T_{\mathrm{hold}}$~control steps.
Each obstacle zone $T_k$ is a cylindrical zone of radius
$r_k \sim \mathcal{U}(r_{\min}^{\mathrm{obs}}, r_{\max}^{\mathrm{obs}})$,
where $r_{\min}^{\mathrm{obs}}$ and $r_{\max}^{\mathrm{obs}}$ are the minimum
and maximum allowable obstacle radii.
Obstacle zones move on the $xy$-plane at speed~$v_{\mathrm{obs}}$ following
the same random-heading protocol as target zones.
Boundary reflections keep all entities within the task region.

\subsubsection{Safety Constraints}

To avoid collisions and remain within the task region, the UAV must maintain
safe distances.
Let $\mathbf{p}_k^{(t)}$ denote the centre of obstacle zone~$T_k$ at
time~$t$.
Let $R_{\mathrm{safe}}^{1} = r_{\mathrm{uav}} + r_k$ denote the minimum safe
distance between the UAV and an obstacle zone boundary, and
$R_{\mathrm{safe}}^{2} = r_{\mathrm{uav}}$ the minimum safe distance
between the UAV and the task region boundary.
The constraints are
\begin{equation}
  \begin{cases}
    \lVert\mathbf{p}^{(t)} - \mathbf{p}_k^{(t)}\rVert
      \ge R_{\mathrm{safe}}^{1},
    & \forall\, T_k \in \mathcal{T}, \\[4pt]
    d_{\mathrm{b}}^{(t)} \ge R_{\mathrm{safe}}^{2},
  \end{cases}
  \label{eq:safety}
\end{equation}
where
$d_{\mathrm{b}}^{(t)} = \min\!\bigl(
  \tfrac{L}{2} - |x^{(t)}|,\;
  \tfrac{L}{2} - |y^{(t)}|
\bigr)$
is the minimum distance from the UAV centre to the task region boundary.

% ==============================================================
\subsection{Problem Transformation}
\label{subsec:problem_transformation}

The target coverage task is modelled as a
\emph{Partially Observable Markov Decision Process}~(POMDP),
defined by the tuple
$\langle \mathcal{S}, \mathcal{O}, \mathcal{A}, \mathcal{P}, R, \gamma
\rangle$,
where $\mathcal{S}$ is the global state space,
$\mathcal{O}$ is the observation space,
$\mathcal{A}$ is the action space,
$\mathcal{P}(\mathbf{s}' \mid \mathbf{s}, a)$ is the state transition function,
$R: \mathcal{S} \times \mathcal{A} \to \mathbb{R}$ is the reward function, and
$\gamma \in (0,1]$ is the discount factor.
The UAV follows a parameterised policy~$\pi_{\theta}(a \mid o)$ with learnable
parameters~$\theta$, aiming to maximise the expected discounted return
$G_t = \sum_{k=0}^{\infty} \gamma^k r_{t+k}$, where $r_{t+k}$ is the reward
received at time $t+k$.

\subsubsection{Observation Space}
\label{subsubsec:observation}

At each time step~$t$, the UAV receives a local observation $o^{(t)}$ composed
of three parts:
\begin{enumerate}
  \item \emph{Self-state}
    $[\mathbf{p}^{(t)},\, \mathbf{v}^{(t)}]$:
    the position and velocity of the UAV, normalised by the region half-extent~$L/2$
    and maximum speed~$v_{\max}$, respectively.
  \item \emph{Target state}
    $[\mathbf{p}_m^{(t)} - \mathbf{p}^{(t)}]$,
    $m = 1,\dots,M$:
    the relative position of each target with respect to the UAV, normalised.
  \item \emph{Obstacle zone state}
    $[\mathbf{p}_k^{(t)} - \mathbf{p}^{(t)},\, r_k]$,
    $k = 1,\dots,K$:
    the relative position and radius of each obstacle zone.
    The radius is normalised by $R_{\mathrm{obs}}$.
\end{enumerate}
For entities beyond the observation radius~$R_{\mathrm{obs}}$, the
corresponding entries are set to the zero vector, indicating no observation.

\subsubsection{Action Space}
\label{subsubsec:action}

The action is a two-dimensional velocity direction on the $xy$-plane,
$a^{(t)} = (a^{x}, a^{y}) \in [-1,1]^{2}$.
The policy~$\pi_{\theta}(a \mid o)$ outputs normalised actions, which are
then scaled by $v_{\max}$ and converted to rotor RPMs by the cascaded PID
controller described in Section~\ref{subsec:system_model}.

\subsubsection{Reward Function}
\label{subsubsec:reward}

The reward function is inspired by the Artificial Potential Field~(APF) method,
combining repulsive penalties that grow as the UAV approaches dangerous regions
with attractive rewards that guide it toward targets.
It decomposes into three components: obstacle zone penalty, target reward,
and boundary penalty.
A common potential kernel is used throughout:
\begin{equation}
  \varphi(d,\,d_0) =
    \alpha\,
    \left(
      \frac{1}{d - d_0 + 1}
      - \frac{1}{\sigma + 1}
    \right)^{\!2},
  \label{eq:apf_kernel}
\end{equation}
where $\alpha$ is a tunable scaling coefficient, $d$ is the current distance,
$d_0$ is the contact distance at which hard collision occurs, and
$\sigma$ is the width of the threat zone surrounding the UAV.
The potential $\varphi$ equals zero at $d = d_0 + \sigma$ and increases
monotonically as $d \to d_0$.
Throughout this section, the subscript~$xy$ denotes the projection of a
three-dimensional position onto the $xy$-plane, e.g.\
$\mathbf{p}_{xy} = (x, y)^{\top}$.

\paragraph{Obstacle zone penalty.}
For obstacle zone~$T_k$ with radius~$r_k$, the planar distance between the
UAV and the obstacle zone centre is
$d = \lVert\mathbf{p}_{xy} - \mathbf{p}_{k,xy}\rVert$,
where $\mathbf{p}_{xy}$ and $\mathbf{p}_{k,xy}$ are the $xy$-projections of
the UAV and obstacle zone positions, respectively.
Let $D_k = r_{\mathrm{uav}} + r_k$ denote the contact distance.
The penalty is
\begin{equation}
  r_{\mathrm{obs}} =
  \begin{cases}
    -20 & d \le D_k, \\[3pt]
    -\varphi(d,\, D_k)
    & D_k < d < D_k + \sigma, \\[3pt]
    0 & \text{otherwise}.
  \end{cases}
  \label{eq:r_obs}
\end{equation}

\paragraph{Target reward.}
Let $d^{*} = \min_{m} \lVert\mathbf{p}_{xy} - \mathbf{p}_{m,xy}\rVert$
denote the planar distance to the nearest target zone centre,
where $\mathbf{p}_{m,xy}$ is the $xy$-projection of target zone~$G_m$.
The reward is
\begin{equation}
  r_{\mathrm{tgt}} =
  \begin{cases}
    +10
    & d^{*} \le R_{\mathrm{cov}}, \\[3pt]
    \varphi(d^{*},\, R_{\mathrm{cov}})
    & R_{\mathrm{cov}} < d^{*} < R_{\mathrm{cov}} + \sigma, \\[3pt]
    0 & \text{otherwise}.
  \end{cases}
  \label{eq:r_tgt}
\end{equation}

\paragraph{Boundary penalty.}
Using $d_{\mathrm{b}}^{(t)}$ defined in~\eqref{eq:safety}, the penalty is
\begin{equation}
  r_{\mathrm{bnd}} =
  \begin{cases}
    -20 & d_{\mathrm{b}} \le r_{\mathrm{uav}}, \\[3pt]
    -\varphi(d_{\mathrm{b}},\, r_{\mathrm{uav}})
    & r_{\mathrm{uav}} < d_{\mathrm{b}} < r_{\mathrm{uav}} + \sigma, \\[3pt]
    0 & \text{otherwise}.
  \end{cases}
  \label{eq:r_bnd}
\end{equation}

The total reward at each step is
\begin{equation}
  R = r_{\mathrm{obs}} + r_{\mathrm{tgt}} + r_{\mathrm{bnd}}.
  \label{eq:total_reward}
\end{equation}

\subsubsection{Episode Termination}
\label{subsubsec:termination}

An episode terminates when:
(i)~the cumulative number of captures reaches
$\lceil \eta \cdot M \rceil$, where $\eta \in (0,1]$ is the coverage
completion ratio;
(ii)~any safety constraint in~\eqref{eq:safety} is violated; or
(iii)~the episode duration exceeds the maximum time horizon~$T_{\max}$.
